%!TEX program = xelatex
% 【關鍵設定】fontset=none 代表不使用預設字型，我們要自己設定
\documentclass[fontset=none, a4paper, 12pt]{ctexart} 

% ==========================================
% 1. 套件載入區
% ==========================================
\usepackage{geometry}
\usepackage{longtable}
\usepackage{fontspec} % 必須載入，用於讀取 ttf
\usepackage{setspace}
\usepackage{titlesec}
\usepackage{titletoc}
\usepackage{tocloft}
\usepackage{fancyhdr}
\usepackage{graphicx}
\usepackage{caption}
% 圖表與標題間距壓縮
\captionsetup[figure]{skip=4pt}
\setlength{\intextsep}{7pt}
\setlength{\textfloatsep}{7pt}
\setlength{\floatsep}{7pt}
\usepackage[hidelinks]{hyperref}
\usepackage{booktabs}
\usepackage{float}
\usepackage{amsmath}
\usepackage{amssymb}
\usepackage{tabularx}
\usepackage{placeins}
\usepackage{needspace}
\usepackage{algorithm}
\usepackage{algpseudocode}
\floatname{algorithm}{演算法}
\usepackage{makecell}
\usepackage{array}
\usepackage{pdflscape}
\usepackage{rotating}
\usepackage{pdfpages} % 插入口試審核頁掃描 PDF
\sloppy


% =====  refs.bib 管理文獻的關鍵設定 =====
\usepackage[
  backend=biber,
  style=numeric,
  sorting=none,
  maxbibnames=99,
  giveninits=true,
  doi=true,
  url=true,
  isbn=false
]{biblatex}

\addbibresource{refs.bib}
% 作者改成：姓在前、名字縮寫在後
\DeclareNameAlias{default}{family-given}

% 最後一位作者前使用 &，比較接近 APA
\DeclareDelimFormat{finalnamedelim}{\addcomma\space\&\space}

% article 不顯示 In:
\renewbibmacro{in:}{%
  \ifentrytype{article}
    {}
    {\printtext{\bibstring{in}\intitlepunct}}}

% 文章標題不要加引號
\DeclareFieldFormat[article]{title}{#1}
\DeclareFieldFormat[inproceedings]{title}{#1}
\DeclareFieldFormat[misc]{title}{#1}

% ==========================================
% 2. 字型設定
% ==========================================

% 整份論文的中文主字型：標楷體
\setCJKmainfont{kaiu.ttf}[
    AutoFakeBold=2.5,
    AutoFakeSlant=0.3
]
% 整份論文的英文字型
\setmainfont{TeX Gyre Termes}
% ------------------------------------------
% 封面專用字型
% ------------------------------------------
% 封面中文：標楷體
\newCJKfontfamily\coverzh{kaiu.ttf}[
    AutoFakeBold=2.5,
    AutoFakeSlant=0.3
]
% 封面英文：Times New Roman
\newfontfamily\coveren{Times New Roman}


% ====================================
% 3. 版面與格式設定
% ====================================

\geometry{a4paper, top=3cm, bottom=3cm, left=3cm, right=3cm}

% 【頁碼規範】頁碼位置應置於每頁的頁尾，對齊方式為置中
\pagestyle{plain} 

% 圖表標題中文化
\captionsetup[figure]{name=圖}
\captionsetup[table]{name=表}
\ctexset{bibname = {參考文獻}}

% 章節標題格式
\titleformat{\section}[block]{\centering\fontsize{20}{24}\selectfont\bfseries}{\thesection}{1em}{}
\titleformat{\subsection}[block]{\fontsize{16}{20}\selectfont\bfseries}{\thesubsection.}{1em}{}
\titleformat{\subsubsection}[block]{\fontsize{14}{18}\selectfont\bfseries}{\thesubsubsection.}{1em}{}

% 編號格式
\renewcommand{\thesection}{第\chinese{section}章}
\renewcommand{\thesubsection}{\arabic{section}.\arabic{subsection}}
\renewcommand{\thesubsubsection}{\arabic{section}.\arabic{subsection}.\arabic{subsubsection}}

% ==========================================
% 目錄、表目錄、圖目錄設定
% ==========================================

% 標題名稱
\renewcommand{\contentsname}{目錄}
\renewcommand{\listtablename}{表目錄}
\renewcommand{\listfigurename}{圖目錄}

% ---------- 目錄標題置中 ----------
\renewcommand{\cfttoctitlefont}{\hfill\Large\bfseries}
\renewcommand{\cftaftertoctitle}{\hfill\null}

% ---------- 表目錄標題置中 ----------
\renewcommand{\cftlottitlefont}{\hfill\Large\bfseries}
\renewcommand{\cftafterlottitle}{\hfill\null}

% ---------- 圖目錄標題置中 ----------
\renewcommand{\cftloftitlefont}{\hfill\Large\bfseries}
\renewcommand{\cftafterloftitle}{\hfill\null}

% ---------- 標題與內容距離 ----------
\setlength{\cftbeforetoctitleskip}{0pt}
\setlength{\cftaftertoctitleskip}{1.2em}

\setlength{\cftbeforelottitleskip}{0pt}
\setlength{\cftafterlottitleskip}{1.2em}

\setlength{\cftbeforeloftitleskip}{0pt}
\setlength{\cftafterloftitleskip}{1.2em}

% ---------- 目錄項目點線 ----------
% 使誌謝辭、摘要、Abstract、目錄、表目錄、圖目錄也顯示 ........ i
\renewcommand{\cftsecleader}{\cftdotfill{\cftdotsep}}
\renewcommand{\cftsubsecleader}{\cftdotfill{\cftdotsep}}
\renewcommand{\cftsubsubsecleader}{\cftdotfill{\cftdotsep}}

% ---------- 目錄字體 ----------
\renewcommand{\cftsecfont}{\bfseries}
\renewcommand{\cftsecpagefont}{\bfseries}

% ---------- 章節間距 ----------
\setlength{\cftbeforesecskip}{1.5em}


\begin{document}


% ==========================================
% 封面（不顯示頁碼）
% ==========================================

% 封面上下左右邊界各 3 公分
\newgeometry{
    a4paper,
    top=3cm,
    bottom=3cm,
    left=3cm,
    right=3cm
}

\begin{titlepage}
    \centering
    \thispagestyle{empty}

    \vspace*{0.7cm}

    % 中文校名與學程名稱：標楷體 20pt，同一行
    {\coverzh
    \fontsize{20pt}{28pt}\selectfont
    \begingroup
    \xeCJKsetup{CJKglue=\hskip -1.55pt}
    \makebox[\textwidth][c]{%
        國立中興大學人工智慧與資料科學碩士在職學位學程%
    }
    \endgroup
    \par}

    \vspace{0.35cm}

    % 論文類別
    {\coverzh
    \fontsize{20pt}{28pt}\selectfont
    碩士學位論文\par}

    % 論文類別與中文題目之間
    \vspace{2.5cm}

    % 中文論文題目：倒三角形排列
    {\coverzh
    \fontsize{20pt}{29pt}\selectfont
    XXXX\par}

    \vspace{0.75cm}

    % 英文論文題目：Times New Roman 20pt
    {\coveren
    \fontsize{20pt}{26pt}\selectfont
    \makebox[\textwidth][c]{%
     
    }\par}

    % 論文題目與指導教授資訊之間
    \vspace{2.4cm}

    % 指導教授
    {\coverzh
    \fontsize{20pt}{28pt}\selectfont
    指導教授：
    \hspace{0.3em}
    林長鋆
    \hspace{0.1em}
    {\coveren Chang-Yun Lin}\par}

    \vspace{0.35cm}

    % 研究生
    {\coverzh
    \fontsize{20pt}{28pt}\selectfont
    研究生：
    \hspace{0.3em}
    XXX
    \hspace{0.1em}
    {\coveren XXX}\par}

    % 研究生與完成日期之間
    \vspace{2.0cm}

    % 論文完成日期
    {\coverzh
    \fontsize{20pt}{28pt}\selectfont
    中華民國一百一十X年八月\par}

\end{titlepage}

% 恢復論文正文原本的邊界設定
\restoregeometry

\clearpage

% ==========================================
% 口試審核頁（不編頁碼）
% ==========================================
% 尚未取得口試審核頁時，保留以下空白頁。
\thispagestyle{empty}
\vspace*{\fill}
\begin{center}
    \Large 口試審核頁預留
\end{center}
\vspace*{\fill}
\clearpage

% ------------------------------------------------------------
% 【口試審核頁 PDF 插入方法】
% 取得掃描 PDF 後：
% 1. 將上面的「口試審核頁預留」區塊刪掉或註解。
% 2. 將 PDF 放入 picture/ 資料夾。
% 3. 取消下方程式碼的註解。
%
% \includepdf[
%     pages=1,
%     fitpaper=false,
%     trim=0mm 1.5mm 1.5mm 0mm,
%     clip,
%     pagecommand={\thispagestyle{empty}}
% ]{picture/口試審核頁檔名.pdf}
% \clearpage
% ------------------------------------------------------------

\clearpage


% ==========================================
% 誌謝（置於中文摘要之前，不顯示頁碼，目錄不用放）
% ==========================================

\clearpage
\pagenumbering{gobble}
\pagestyle{empty}

\begin{center}
    \large\bfseries 誌謝
\end{center}
\vspace{0.5em}

.........


\vspace{2em}

\begin{flushright}
XXX　謹誌\\
中華民國一百一十X年八月
\end{flushright}


%======================================
% 前半部：羅馬數字編碼 (i, ii, iii...)
% 規範：摘要至正文前，以小寫羅馬數字連續編碼
% =====================================

\clearpage
\pagestyle{plain}
\pagenumbering{roman}
\setcounter{page}{1}

% ==========================================
% 中文摘要
% ==========================================
\phantomsection
\addcontentsline{toc}{section}{摘要}
\begin{center}
\large\bfseries 摘要
\end{center}
\vspace{0.5em}

.....

\vspace{1em}
\noindent \textbf{關鍵詞：}

\newpage

% ==========================================
% 英文摘要 (Abstract)
% ==========================================
\phantomsection
\addcontentsline{toc}{section}{Abstract}
\begin{center}
\large\bfseries Abstract
\end{center}
\vspace{0.5em}

......

\vspace{1em}
\noindent \textbf{Keywords:}

\newpage


% ==========================================
% 目錄
% ==========================================
\phantomsection
\addcontentsline{toc}{section}{目錄}
\tableofcontents
\newpage

% 表目錄
\phantomsection
\addcontentsline{toc}{section}{表目錄}
\listoftables
\newpage

% 圖目錄
\phantomsection
\addcontentsline{toc}{section}{圖目錄}
\listoffigures
\newpage

% ==========================================
% 正文：阿拉伯數字頁碼 1, 2, 3...
% ==========================================
\clearpage
\pagenumbering{arabic}
\setcounter{page}{1}
\begin{doublespace}

% ============================================================
% 第一章 緒論
% ============================================================
\section{緒論}

\subsection{研究背景與動機}
% 請撰寫內容。

\subsection{研究目的}
% 請撰寫內容。

% 如需三級標題，可使用：
% \subsubsection{次小節標題}
% 請撰寫內容。

% ============================================================
% 第二章 文獻探討
% ============================================================
\clearpage
\section{文獻探討}

\subsection{小節標題}
% 請依研究主題自行修改小節名稱與內容。

\subsubsection{次小節標題}
% 請撰寫內容。

% ============================================================
% 圖片插入範例（一般直式圖片）
% ============================================================
% 使用方式：將圖片放在 picture/，取消以下註解並修改檔名、圖名、label。
%
% 文中引用寫法：如圖~\ref{fig:example} 所示。
%
% \begin{figure}[H]
%     \centering
%     \includegraphics[
%         width=0.82\textwidth,
%         keepaspectratio
%     ]{picture/your_figure.png}
%     \caption{請輸入圖片標題}
%     \label{fig:example}
% \end{figure}
%
% 寬度常用值：0.7\textwidth、0.82\textwidth、0.95\textwidth。
% 若圖片過高，可加入 height=0.55\textheight。

% ============================================================
% 表格插入範例（tabularx，自動分配欄寬）
% ============================================================
% 文中引用寫法：如表~\ref{tab:example} 所示。
%
% \begin{table}[H]
% \centering
% \caption{請輸入表格標題}
% \label{tab:example}
% \footnotesize
% \renewcommand{\arraystretch}{1.12}
% \setlength{\tabcolsep}{4pt}
% \begin{tabularx}{0.96\textwidth}{
% >{\raggedright\arraybackslash}p{3.0cm}
% >{\centering\arraybackslash}p{2.0cm}
% >{\raggedright\arraybackslash}X}
% \hline
% 欄位一 & 欄位二 & 欄位三 \\
% \hline
% 範例 A & 1 & 請填寫內容 \\
% 範例 B & 2 & 請填寫內容 \\
% \hline
% \end{tabularx}
% \end{table}

% ============================================================
% 第三章 研究方法
% ============================================================
\clearpage
\section{研究方法}

\subsection{研究架構}
% 請撰寫內容。

\subsection{資料來源與研究設計}
% 請撰寫內容。

\subsection{研究方法或模型}
% 請撰寫內容。

% ============================================================
% 橫式大型圖片插入範例（只旋轉圖片與圖說）
% ============================================================
% 適合流程圖、模型架構圖、寬版圖。
%
% \clearpage
% \begin{sidewaysfigure}[p]
%     \centering
%     \includegraphics[
%         width=1.05\textheight,
%         height=0.98\textwidth,
%         keepaspectratio
%     ]{picture/your_wide_figure.png}
%     \caption{請輸入橫式圖片標題}
%     \label{fig:wide-example}
% \end{sidewaysfigure}
% \clearpage

% ============================================================
% 同一圖放兩張圖片範例
% ============================================================
% \begin{figure}[H]
% \centering
% \begin{minipage}{0.92\textwidth}
% \centering
%
% \includegraphics[width=\linewidth,height=0.32\textheight,keepaspectratio]
% {picture/figure_a.png}
% \vspace{0.25em}
% {\scriptsize (a) 子圖 A}
%
% \vspace{0.8em}
%
% \includegraphics[width=\linewidth,height=0.32\textheight,keepaspectratio]
% {picture/figure_b.png}
% \vspace{0.25em}
% {\scriptsize (b) 子圖 B}
%
% \end{minipage}
% \caption{請輸入整組圖片標題}
% \label{fig:two-images}
% \end{figure}

% ============================================================
% 公式插入範例（如不需要可直接刪除此註解區）
% ============================================================
% 文中引用：由式~\eqref{eq:example} 可知……
%
% \begin{equation}
% y = ax + b
% \label{eq:example}
% \end{equation}

% ============================================================
% 第四章 實驗結果與討論
% ============================================================
\clearpage
\section{實驗結果與討論}

\subsection{實驗結果}
% 請撰寫內容。

\subsection{結果分析與討論}
% 請撰寫內容。

% ============================================================
% 第五章 結論與未來展望
% ============================================================
\clearpage
\section{結論與未來展望}

\subsection{研究成果概述}
% 請撰寫內容。

\subsection{研究限制}
% 請撰寫內容。

\subsection{未來研究方向}
% 請撰寫內容。

\subsection{總結}
% 請撰寫內容。

% ============================================================
% 文獻引用方法
% ============================================================
% refs.bib 中先建立文獻，例如 example2026。
% 正文可寫：作者提出相關方法~\cite{example2026}。
% 多篇引用：\cite{example2026,exampleBook2025}

% ==========================================
% 參考文獻
% ==========================================
\clearpage
\phantomsection
\printbibliography[heading=bibintoc,title={參考文獻}]

\end{doublespace}
\end{document}
